\section{Archived hardware confidence audit}
\label{app:hardware-confidence}
The archived data comprise 152 W-state runs and 15 higher-excitation Dicke
runs. The 30 tabulated results match the archived excitation numbers,
system sizes, energies, and execution dates. The runtime metadata of all
152 W-state jobs record measurement-error mitigation and measurement
twirling as enabled; those of all 15 higher-excitation Dicke jobs record
both as disabled. Gate twirling, dynamical decoupling, zero-noise
extrapolation, and probabilistic error cancellation are disabled in all
167 jobs. The configured shot counts are 16384 (W, divided among 32
measurement-twirling randomizations) and 16000 (higher excitation, one
randomization); the runtime options apply this count to each
measurement-basis configuration. The archived results record these
configured options but not the executed basis-resolved circuits or
outcomes. The W-state results
select the lowest measured energy at each size. Figure~\ref{fig:W_experiment}
shows the runtime-reported uncertainty of the selected result.
For the unmitigated results we evaluate the range-only bound in
Appendix~\ref{app:resources}. With $N_b=N_0$, let $R_b$ be the range of
the per-shot basis score. Conditional on independent preparations,
trusted ideal Pauli measurements, and $N_0$ shots in each of the three
separately sampled bases, define
\begin{equation}
 E_U=\widehat E+
 \sqrt{\frac{R_X^2+R_Y^2+R_Z^2}{2N_0}\ln\frac{167}{0.05}}.
\end{equation}
For each unmitigated result, Hoeffding's inequality then gives
$\Pr\{E>E_U\}\le0.05/167$ for its true energy $E$. A union bound therefore gives simultaneous coverage
of at least $95\%$ for all 15 higher-excitation results. The denominator
167 includes the full archived collection and accounts for selection
within it. The union bound permits dependence between runs. Its coverage
is restricted to the archived collection and the stated sampling model;
selection from additional trials would require a corresponding adjustment.
The calculation takes the configured shot allocation, $N_0=16000$ per
homogeneous basis, as the executed sampling model, which the archived
records do not independently confirm, and assumes ideal readout. Assessing readout bias or mitigated W-state results requires
an error model beyond these assumptions.

For the four passing higher-excitation runs the conditional bounds are:
\begin{center}
\begin{tabular}{ccccc}
$n$ & $k$ & $\widehat E$ & $E_U$ & $1-\alpha$\\\hline
4 & 2 & 0.136078 & 0.188890 & 0.333333\\
5 & 2 & 0.138325 & 0.265433 & 0.400000\\
4 & 3 & 0.078969 & 0.212253 & 0.250000\\
5 & 3 & 0.151250 & 0.278358 & 0.400000\\
\end{tabular}
\end{center}
Both higher-excitation families also retain positive conditional fidelity
bounds $\max\{0,1-E_U\}$ through eight qubits. No larger archived member
passes this range-only positivity test. A failed test leaves entanglement undetermined. At seven qubits the $k=2,3$ upper bounds are respectively
$0.623486$ and $0.528390$, exceeding their witness thresholds $2/7$ and
$3/7$. Thus the seven-qubit point-estimate crossings are not promoted to
finite-confidence GME claims here.

The available data contain aggregate energy uncertainties. The confidence
bounds above use the recorded shot counts and the possible ranges of the
measurement scores. Basis-resolved outcome statistics would be needed to
evaluate their empirical variances. The six- and seven-qubit
higher-excitation crossings and all W-state crossings are therefore
reported as point estimates under the present analysis.
The audit script, all-run and table-matched CSV files, and input hashes
are supplied in the revision's source supplement under \texttt{pra\_confidence\_20261013/analysis/} and \texttt{pra\_confidence\_20261013/audit-results/}.
